% Formal Q26 validation-only spectral-CNN ablation; paired E130/E190 checkpoints within each seed.
\begin{table}[htbp]
\centering\scriptsize
\setlength{\tabcolsep}{2.4pt}
\bicaption[tab:component-results]{表}{正式频谱CNN消融的逐seed绝对误差（44名验证被试，Q26）}{Tab.}{Per-seed absolute errors in the formal spectral-CNN ablation (44 validation subjects, Q26)}
\resizebox{\linewidth}{!}{%
\begin{tabular}{lrrrrrr}
\toprule
& \multicolumn{2}{c}{Seed 20260821} & \multicolumn{2}{c}{Seed 20260822} & \multicolumn{2}{c}{Seed 20260823} \\
指标$\downarrow$ & E130 & E190 & E130 & E190 & E130 & E190 \\
\midrule
全域ERB（dB） & $0.8554\pm0.1635$ & $0.8304\pm0.1645$ & $0.8501\pm0.1599$ & $0.8257\pm0.1610$ & $0.8488\pm0.1607$ & $0.8277\pm0.1620$ \\
对侧$25^\circ$ ERB（dB） & $1.2600\pm0.1911$ & $1.2433\pm0.1918$ & $1.2616\pm0.1882$ & $1.2407\pm0.1881$ & $1.2540\pm0.1881$ & $1.2397\pm0.1909$ \\
频带ILD（dB） & $1.6132\pm0.1659$ & $1.5374\pm0.1541$ & $1.6180\pm0.1705$ & $1.5389\pm0.1555$ & $1.6105\pm0.1761$ & $1.5372\pm0.1631$ \\
ITD加权MAE（$\mu$s） & $15.5683\pm2.3349$ & $15.5411\pm2.3330$ & $15.5621\pm2.3373$ & $15.5520\pm2.3399$ & $15.5431\pm2.3281$ & $15.5147\pm2.3260$ \\
全域LSD（dB） & $3.8485\pm0.2896$ & $3.5968\pm0.2692$ & $3.8352\pm0.2937$ & $3.5923\pm0.2773$ & $3.8471\pm0.2853$ & $3.5896\pm0.2672$ \\
高频一阶差分（dB/bin） & $0.4930\pm0.0738$ & $0.4128\pm0.0549$ & $0.4891\pm0.0710$ & $0.4110\pm0.0542$ & $0.4878\pm0.0691$ & $0.4122\pm0.0542$ \\
高频二阶差分（dB/bin$^2$） & $0.3082\pm0.0694$ & $0.2360\pm0.0488$ & $0.3017\pm0.0653$ & $0.2328\pm0.0473$ & $0.3006\pm0.0614$ & $0.2356\pm0.0469$ \\
多尺度凹口深度（dB） & $0.5667\pm0.0837$ & $0.4921\pm0.0660$ & $0.5640\pm0.0822$ & $0.4912\pm0.0656$ & $0.5630\pm0.0819$ & $0.4917\pm0.0657$ \\
\bottomrule
\end{tabular}%
}
\par\smallskip
\parbox{\linewidth}{\footnotesize 数值为被试均值$\pm$样本标准差。E130为不含频谱CNN的FiLM-SIREN末点，E190为同seed E130基础上的完整FSC末点；所有checkpoint、epoch、seed和超参数均预先冻结。}
\end{table}

\begin{table}[htbp]
\centering\scriptsize
\setlength{\tabcolsep}{3pt}
\bicaption[tab:component-paired]{表}{正式频谱CNN消融的逐seed配对差及95\% bootstrap区间}{Tab.}{Per-seed paired differences and 95\% bootstrap intervals in the formal spectral-CNN ablation}
\resizebox{\linewidth}{!}{%
\begin{tabular}{lrrr}
\toprule
指标$\downarrow$ & Seed 20260821 & Seed 20260822 & Seed 20260823 \\
\midrule
全域ERB（dB） & $-0.0250\;[-0.0275,-0.0225]$ & $-0.0244\;[-0.0271,-0.0218]$ & $-0.0210\;[-0.0235,-0.0186]$ \\
对侧$25^\circ$ ERB（dB） & $-0.0167\;[-0.0199,-0.0136]$ & $-0.0208\;[-0.0250,-0.0169]$ & $-0.0142\;[-0.0179,-0.0105]$ \\
频带ILD（dB） & $-0.0758\;[-0.0854,-0.0661]$ & $-0.0791\;[-0.0889,-0.0695]$ & $-0.0733\;[-0.0831,-0.0638]$ \\
ITD加权MAE（$\mu$s） & $-0.0271\;[-0.0378,-0.0170]$ & $-0.0101\;[-0.0199,-0.0007]$ & $-0.0284\;[-0.0388,-0.0182]$ \\
全域LSD（dB） & $-0.2517\;[-0.2670,-0.2369]$ & $-0.2429\;[-0.2574,-0.2289]$ & $-0.2576\;[-0.2722,-0.2434]$ \\
高频一阶差分（dB/bin） & $-0.0802\;[-0.0868,-0.0742]$ & $-0.0781\;[-0.0841,-0.0727]$ & $-0.0757\;[-0.0811,-0.0707]$ \\
高频二阶差分（dB/bin$^2$） & $-0.0722\;[-0.0789,-0.0659]$ & $-0.0689\;[-0.0748,-0.0633]$ & $-0.0649\;[-0.0697,-0.0602]$ \\
多尺度凹口深度（dB） & $-0.0746\;[-0.0818,-0.0683]$ & $-0.0728\;[-0.0795,-0.0669]$ & $-0.0713\;[-0.0779,-0.0656]$ \\
\bottomrule
\end{tabular}%
}
\par\smallskip
\parbox{\linewidth}{\footnotesize 差值定义为$\Delta=E_{\mathrm{FSC\ E190}}-E_{\mathrm{FiLM\text{-}SIREN\ E130}}$，负值表示加入频谱CNN refinement后误差降低。每个seed分别对相同44名验证被试进行10000次paired bootstrap（固定bootstrap seed 20260909）；三个seed不合并为132个独立样本。}
\end{table}
